% Résultats RPEM sur machine, découpés en trois parties (à utiliser dans une frame avec \only<1-3>)
%   partie 1 : V_c + a_w, b_w, c_w
%   partie 2 : V_c + K_f, K_c, K_m et les efforts associés
%   partie 3 : V_c + W_c,max et critère J
% Les sous-figures sont découpées dans images/rpem.pdf_tex (grille 3 lignes x 4 colonnes).
% Filtre vert = résultat validé (tout sauf K_m et M_P,ez).
% Réglage : \rpemlargeur = largeur de la figure MATLAB complète.
% La macro \rpempanneau est définie dans le préambule de presentation.tex.
\providecommand{\rpemlargeur}{0.9\linewidth}
\sbox{\rpemboite}{\fontsize{3.5pt}{5pt}\selectfont\def\svgwidth{\rpemlargeur}\import{images/}{rpem.pdf_tex}}%
\setlength{\rpemL}{\rpemlargeur}% largeur de la figure, pour les épaisseurs de trait redessinées
\begin{tikzpicture}[x=\rpemlargeur, y=\rpemlargeur]
    % bornes des cellules dans la figure d'origine (fraction de la largeur)
    \def\cA{0}\def\cB{0.251}\def\cC{0.500}\def\cD{0.745}\def\cE{1}
    \def\rA{0}\def\rB{0.207}\def\rC{0.417}\def\rD{0.626}
    % disposition : V_c agrandi à gauche, panneaux de chaque partie à droite (2 lignes max)
    \def\pas{0.258}\def\haut{0.224}% pas des colonnes, décalage entre les deux lignes
    % V_c agrandi pour occuper exactement la hauteur des deux lignes de droite
    \pgfmathsetmacro{\ech}{(\haut+\rD-\rC)/(\rD-\rC)}
    \def\yvc{0}
    \pgfmathsetmacro{\xg}{\ech*(\cB-\cA)+0.012}

    % V_c : présent dans les trois parties
    \rpempanneau[echelle=\ech]{\cA}{\cB}{\rC}{\rD}{0}{\yvc}{1}

    % partie 1 : coefficients du modèle a_w, b_w, c_w
    \only<1>{%
        % fenêtre remontée de 0,006 : les exposants (x10^n) dépassent le haut de la figure MATLAB
        \pgfmathsetmacro{\rCh}{\rC+0.006}\pgfmathsetmacro{\rDh}{\rD+0.006}%
        \rpempanneau[droite=0.006, droitey=0.03]{\cB}{\cC}{\rCh}{\rDh}{\xg}{\haut}{1}                    % a_w (sans le bout d'étiquette de b_w)
        \rpempanneau{0.494}{\cD}{\rCh}{\rDh}{\xg+\pas}{\haut}{1}                      % b_w (élargi à gauche : haut du « b » de l'étiquette)
        \rpempanneau[expo-]{\cD}{\cE}{\rCh}{\rDh}{\xg}{0}{1}}         % c_w

    % partie 2 : pressions spécifiques (ligne 1) et efforts associés (ligne 2)
    \only<2>{%
        \rpempanneau[etiq-]{\cA}{\cB}{\rB}{\rC}{\xg}{\haut}{1}          % K_f
        \rpempanneau[etiq-]{\cB}{\cC}{\rB}{\rC}{\xg+\pas}{\haut}{1}     % K_c
        \rpempanneau[etiq-]{\cC}{\cD}{\rB}{\rC}{\xg+2*\pas}{\haut}{0}   % K_m (non exploitable)
        \rpempanneau[etiq+]{\cA}{\cB}{\rA}{\rB}{\xg}{0}{1}              % F_er
        \rpempanneau[etiq+]{\cB}{\cC}{\rA}{\rB}{\xg+\pas}{0}{1}         % F_etheta
        \rpempanneau[bande=0.0738, etiq={$M_{P,\vec{e}_z,\text{max}}$ (N.m)}, etiqx=0.029,
            % graduations d'origine masquées avec l'axe qu'elles chevauchaient, puis axe, grille,
            % petits traits, graduations et « 0 » de l'axe des abscisses redessinés (positions issues du SVG)
            apres={\draw[color={rgb,255:red,33;green,33;blue,33}, line width=0.0007\rpemL] (0.0705,0.0334) -- (0.0705,0.1993);
                \draw[color={rgb,255:red,33;green,33;blue,33}, opacity=0.149, line width=0.0007\rpemL] (0.0705,0.1955) -- (0.0742,0.1955);
                \draw[color={rgb,255:red,33;green,33;blue,33}, line width=0.0007\rpemL] (0.0705,0.1955) -- (0.0722,0.1955);
                \rpemgrad[3.5pt]{0.068}{0.1907}{10}
                \draw[color={rgb,255:red,33;green,33;blue,33}, opacity=0.149, line width=0.0007\rpemL] (0.0705,0.1685) -- (0.0742,0.1685);
                \draw[color={rgb,255:red,33;green,33;blue,33}, line width=0.0007\rpemL] (0.0705,0.1685) -- (0.0722,0.1685);
                \rpemgrad[3.5pt]{0.068}{0.1635}{5}
                \draw[color={rgb,255:red,33;green,33;blue,33}, opacity=0.149, line width=0.0007\rpemL] (0.0705,0.1415) -- (0.0742,0.1415);
                \draw[color={rgb,255:red,33;green,33;blue,33}, line width=0.0007\rpemL] (0.0705,0.1415) -- (0.0722,0.1415);
                \rpemgrad[3.5pt]{0.068}{0.1355}{0}
                \draw[color={rgb,255:red,33;green,33;blue,33}, opacity=0.149, line width=0.0007\rpemL] (0.0705,0.1144) -- (0.0742,0.1144);
                \draw[color={rgb,255:red,33;green,33;blue,33}, line width=0.0007\rpemL] (0.0705,0.1144) -- (0.0722,0.1144);
                \rpemgrad[3.5pt]{0.068}{0.1094}{$-5$}
                \draw[color={rgb,255:red,33;green,33;blue,33}, opacity=0.149, line width=0.0007\rpemL] (0.0705,0.0874) -- (0.0742,0.0874);
                \draw[color={rgb,255:red,33;green,33;blue,33}, line width=0.0007\rpemL] (0.0705,0.0874) -- (0.0722,0.0874);
                \rpemgrad[3.5pt]{0.068}{0.0831}{$-10$}
                \draw[color={rgb,255:red,33;green,33;blue,33}, opacity=0.149, line width=0.0007\rpemL] (0.0705,0.0604) -- (0.0742,0.0604);
                \draw[color={rgb,255:red,33;green,33;blue,33}, line width=0.0007\rpemL] (0.0705,0.0604) -- (0.0722,0.0604);
                \rpemgrad[3.5pt]{0.068}{0.0557}{$-15$}
                \draw[color={rgb,255:red,33;green,33;blue,33}, line width=0.0007\rpemL] (0.0705,0.0334) -- (0.0742,0.0334);
                \rpemgrad[3.5pt]{0.068}{0.0286}{$-20$}
                \node[anchor=base, inner sep=0, font=\fontsize{3.5pt}{3.5pt}\selectfont, text=jvtexte] at (0.0705,0.0186) {0};}]
            {\cC}{\cD}{\rA}{\rB}{\xg+2*\pas}{0}{0}       % M_P,ez (non exploitable, graduations réalignées à droite)
        \path (0,0);}% la partie 2 occupe les deux lignes

    % partie 3 : énergie spécifique maximale et critère J
    \only<3>{%
        \rpempanneau[expo+, expo-]{\cD}{\cE}{\rB}{\rC}{\xg}{\haut}{1}      % W_c,max
        \rpempanneau[expo+]{\cD}{\cE}{\rA}{\rB}{\xg}{0}{1}}                % J

    % largeur constante d'une partie à l'autre (la figure ne se décale pas)
    \path (0,0) rectangle (\xg+3*\pas, \haut+\rD-\rC);
    % légende des filtres, sous V_c
    \node[anchor=north west, font=\tiny, text=jvtexte, inner sep=0pt, align=left]
        at (0, \yvc-0.012)
        {\rpemlegende{1}{0}{1}};
\end{tikzpicture}
